\documentclass[11pt,a4paper,landscape]{article}
\usepackage[utf8]{inputenc}
\usepackage[T1]{fontenc}
\usepackage[czech]{babel}
\usepackage{lmodern}
\usepackage[margin=1.2cm]{geometry}
\usepackage{tikz}
\usetikzlibrary{shapes.geometric,trees}
\pagestyle{empty}

\begin{document}
\begin{center}
\begin{tikzpicture}[
  grow=right,
  level 1/.style={level distance=3.6cm, sibling distance=2.5cm},
  level 2/.style={level distance=3.6cm},
  edge from parent/.style={draw=blue!45!black, thick},
  center/.style={
    ellipse,
    draw=blue!60!black,
    fill=blue!12,
    very thick,
    minimum width=5.2cm,
    minimum height=2.4cm,
    align=center,
    font=\Large\bfseries
  },
  branch/.style={
    rectangle,
    rounded corners=4pt,
    draw=blue!55!black,
    fill=blue!8,
    thick,
    minimum width=4.2cm,
    minimum height=1.4cm,
    align=center,
    font=\small\bfseries
  },
  detail/.style={
    rectangle,
    rounded corners=2pt,
    draw=gray!55,
    fill=white,
    text width=8.2cm,
    align=left,
    inner sep=3pt,
    font=\tiny
  }
]

\node[center] {Technologie}
  child { node[branch] {Výpočetní systémy\\{\footnotesize PB151}}
    child { node[detail] {
      Číselné soustavy, vztahy mezi soustavami,\\
      zobrazení celého čísla, aritmetika\\
      Kódy: vnitřní, vnější, detekční a opravné\\
      Procesory: parametry a architektury\\
      Vnitřní a vnější paměti, principy funkce\\
      V/V zařízení a jejich připojování
    } }
  }
  child { node[branch] {Operační systémy\\{\footnotesize PB152, PV004}}
    child { node[detail] {
      Architektura OS a jádra, režimy procesoru\\
      Programovací rozhraní, knihovny\\
      Uživatel, přístupová práva, virtualizace\\
      Virtuální paměť, proces, stránkové tabulky\\
      Vlákna, plánování, souběžnost, uváznutí\\
      POSIX, vznik procesu, copy-on-write
    } }
  }
  child { node[branch] {Souborové systémy\\{\footnotesize PB152, PV004}}
    child { node[detail] {
      Blokové zařízení, bloková vrstva\\
      I/O plánovač, RAID, šifrování disku\\
      Soubory, alokace volného místa, fragmentace\\
      Adresářová struktura, reprezentace na disku\\
      Vstup a výstup mapovaný do paměti
    } }
  }
  child { node[branch] {Sítě a síťové aplikace\\{\footnotesize PB156, PB156cv}}
    child { node[detail] {
      ISO/OSI, TCP/IP, vrstvy, adresace\\
      Fyzická vrstva, signály, kódování, přístup k médiu\\
      Propojování sítí, přepínání, směrování, multicast\\
      Transportní protokoly, zajištěný přenos, spojení\\
      Aplikační protokoly: pošta, soubory, web, DNS\\
      QoS, multimédia, zabezpečení komunikace
    } }
  }
  child { node[branch] {Principy programování\\{\footnotesize PB071, PV080}}
    child { node[detail] {
      Strukturované programování, imperativní jazyk\\
      Paměťový model, správa paměti, dynamická alokace\\
      Uživatelské datové struktury\\
      Nízkoúrovňová práce s pamětí, ukazatele, pole\\
      Ladění (debugging), bezpečné programování
    } }
  }
  child { node[branch] {Databáze\\{\footnotesize PB154 || PB168}}
    child { node[detail] {
      Relační model, schéma, klíče, integritní omezení\\
      Relační algebra, spojování relací\\
      SQL: select, spojování, agregační funkce\\
      Zpracování dotazů, indexování\\
      Transakce, vlastnosti transakčního zpracování
    } }
  }
  child { node[branch] {Softwarové a informační systémy\\{\footnotesize PB007, PV028}}
    child { node[detail] {
      Vývoj a provoz softwarových systémů\\
      UML při vývoji software\\
      Aplikované informační systémy\\
      Architektury rozsáhlých informačních systémů\\
      Metody vedení projektů, řízení provozu IS
    } }
  }
;

\end{tikzpicture}
\end{center}
\end{document}
